\section{Scope, field of definition, and limitations}\label{sec:limitations}

The theorem is an exact planar statement.  It preserves full Boolean atom
data in a polygonal carrier and finitely many one-dimensional sections.  It
does not address polygonal approximation in a metric norm, optimize the
number of vertices, or provide a higher-dimensional filling.

\paragraph{Field of definition.}
Prescribed trace locations are not moved.  Rational finite data lead to a
first-order feasibility problem with rational coefficients and therefore to
an algebraic output, but the proof does not guarantee rational vertices in
general.  Algebraic data yield algebraic output.  If a prescribed endpoint
is transcendental, it must remain so; the construction is interpreted in the
real closure of the ordered field generated by the retained source data.
No universal rationalization claim is made.

\paragraph{Relation to the higher-dimensional conjecture.}
Polytope-convex neural codes are characterized by representable oriented
matroids \citep{KuninLienkaemperRosen2023}.  It remains unknown whether every
open convex code in higher dimension is polytope convex.  The present result
provides exact face and interface data useful in attempts at higher-dimensional
gluing, but it does not settle that conjecture.

\paragraph{Normal data.}
Pointwise equality of edge membership traces is weaker than equality of
finite concave defining functions on the edge, and much weaker than
compatibility of their normal slopes.  The tetrahedral-boundary theorem must
not be read as a convex filling theorem.
